\section{Cloud nach NIST}
Die NIST definiert Cloud nicht über Marketingbegriffe, sondern über 5 technische und organisatorische Eigenschaften.
\subsection{NIST – National Institute of Standards and Technology}
NIST ist eine US-amerikanische Behörde. Die NIST ist gleichzeitig eine Großforschungseinrichtung mit den Forschungsbereichen Messtechnik von physikalisch-technischen Konstanten und heutzutage auch Informationssicherheit.
\subsection{On-Demand Self-Service}
Ein Cloud Nutzer kann eigenständig Services erwerben. Ohne IT-Abteilung, Callcenter oder Vermittler. Auf einer Cloud sollte man jederzeit zugreifen können sowie die Cloud jederzeit kündigen können. Zusätzlich läuft die Cloud über automatisierte Schnittstellen.
\subsection{Broad Network Access}
Eine Cloud sollte von überall mit internetfähigem Gerät erreichbar sein. Da der Cloud-Service flächendeckend verfügbar ist.
\subsection{Resource Pooling}
Bei einer Cloud werden Ressourcen geteilt. Das bedeutet mehrere Nutzer teilen sich die Rechenleistung eines Servers. Um Datenschutz der jeweiligen Region oder Zone zu gewährleisten muss der Cloud Provider entsprechende Maßnahmen treffen.
\subsection{Rapid Elasticity}
Elastizität im Cloud Kontext bedeutet das eine Anwendung automatisch skaliert werden kann, also das ein richtiges Ressourcenlevel aus Rechenleistung, Speicher, Bandbreite und Speicher variabel gewählt wird. Dies ist wichtig um möglichst gut auf Nutzung reagieren zu können ohne das ein Ausfall passiert. Ein beliebtes Beispiel ist hier der Online Shop an Black Friday.
\subsection{Measured Service}
Ein Service in der Cloud sollte messbar sein. Das ist gerade wichtig für Pay-as-you-go, also das man nur das zahlt was man auch wirklich nutzt. Dadurch ist eine höhere Transparenz möglich, wodurch man Kosten besser kalkulieren kann.
